% Section 11.4: Pushforward and pullback | Part: 2 | Depends on: sec:manifolds, sec:vectors, sec:ext | Status: done | Original lines: 1969-2011
%!TEX root = ../../main.tex
\subsection{Pushforward and pullback}\label{sec:pushpull}
% REWRITTEN: pull-back written \Phi\pb instead of \Phi^{*} (same rendering; docs/NOTATION.md)
\begin{formal}[{differential of a map and pullback}]
For a smooth $\Phi:M\to N$ the \emph{differential} (pushforward) is the linear map $T_p\Phi:T_pM\to T_{\Phi(p)}N$, $T_p\Phi(D)(\tilde f)=D(\tilde f\circ\Phi)$; in coordinates it is the Jacobian matrix. If $\Phi$ is a diffeomorphism and $\tilde\tau\in\TT{0}{k}(T_{\Phi(p)}N)$ is a covariant tensor, its \emph{pullback} is
\[
(\Phi\pb\tilde\tau)(v_1,\dots,v_k)=\tilde\tau\big(T_p\Phi(v_1),\dots,T_p\Phi(v_k)\big).
\]
For a general $(\ell,k)$ tensor, covectors are fed in via $\alpha_j\circ T\Phi^{-1}$. Properties: $(\Phi\circ\Psi)\pb=\Psi\pb\circ\Phi\pb$.
\end{formal}
\begin{intuition}[{pullback = ``transport the measuring device backwards''}]
To evaluate a ruler that lives on $N$ on vectors that live on $M$, first push the vectors forward to $N$. Pulling back a metric (the dot product on $N$) gives the \emph{induced metric}: the way one measures lengths of vectors on a surface inside a bigger space. Pulling back with a coordinate chart gives the \emph{components in coordinates}.
\end{intuition}

% REORDERED: the general definition of the pull-back of a form (eq:pullback) was given in Section sec:int1def
% (original L2475-2478); it is merged here with the definition above. The naturality of the wedge (former item (e)
% of prop:wedgeprops, original L987 and its proof L1036) is stated and proved here.
\emph{Forms can be pulled back by any smooth map.} For covariant tensors, and in particular for forms, the formula above
needs no inverse, so $\Phi$ need not be a diffeomorphism. For a smooth map $\phi$ (between manifolds of any dimensions;
$\dd\phi_p=T_p\phi$ is its differential), the pull-back of a form $\omega$ is always defined by contraction with pushed-forward vectors:
\begin{equation}\label{eq:pullback}
(\phi\pb\omega)_p(v)=\omega_{\phi(p)}\bigl(\dd\phi_p\,v\bigr),
\end{equation}
and in the same way, slot by slot, for $k$-forms. On functions ($0$-forms) $\phi\pb f=f\circ\phi$. The pull-back is
\emph{natural} with respect to the wedge product:
\[
\phi\pb(\alpha\wedge\beta)=\phi\pb\alpha\wedge\phi\pb\beta .
\]
\begin{proof}
$(\phi\pb\alpha)(v_1,\dots)=\alpha(\dd\phi\,v_1,\dots)$, and \eqref{eq:wedge} only uses evaluation of $\alpha,\beta$ on
vectors: evaluating both sides on $(v_1,\dots,v_{k+l})$ gives the same sum over permutations, with every $v_i$ replaced by $\dd\phi\,v_i$.
\end{proof}
Together with $\phi\pb\dd=\dd\,\phi\pb$ (Theorem~\ref{thm:dprops}(2) in Section~\ref{sec:dform}) this lets one pull back any form written in coordinates by
substitution, as in the next example. (Step-by-step proofs of naturality, of $(\Phi\circ\Psi)\pb=\Psi\pb\circ\Phi\pb$ and of
$\phi\pb\dd=\dd\,\phi\pb$: Identities~\ref{id:pbwedge}, \ref{id:pbcomp} and~\ref{id:pbd} in the Glossary.)% REWRITTEN: glossary pointer

\begin{example}[{Polar metric}]\label{ex:polarmetric}
The Euclidean metric on $\R^2$ is $g=\dd x\otimes\dd x+\dd y\otimes\dd y$. Substituting $\dd x,\dd y$ from Example~\ref{ex:polarbasis}:
\[
g=(\cos\varphi\,\dd r-r\sin\varphi\,\dd\varphi)^{\otimes2}+(\sin\varphi\,\dd r+r\cos\varphi\,\dd\varphi)^{\otimes2}
=\dd r\otimes\dd r+r^2\,\dd\varphi\otimes\dd\varphi ,
\]
because the mixed terms $\pm r\sin\varphi\cos\varphi\,(\dd r\otimes\dd\varphi+\dd\varphi\otimes\dd r)$ cancel. So $g_{rr}=1$, $g_{\varphi\varphi}=r^2$, $g_{r\varphi}=0$. Consistently, $g(\partial_\varphi,\partial_\varphi)=r^2$.
\end{example}

% REORDERED (F2): moved from Section sec:metric (original L826-831), where the polar metric and the coordinate
% fields were not yet available.
\begin{example}[{raising an index in polar coordinates: the gradient}]
In $\R^2$ with polar coordinates $(r,\varphi)$ the metric components are (Example~\ref{ex:polarmetric}) $g=\mathrm{diag}(1,r^2)$, $g^{-1}=\mathrm{diag}(1,r^{-2})$. The covector $\dd f=\partial_rf\,\dd r+\partial_\varphi f\,\dd\varphi$ becomes, after raising the index (Section~\ref{sec:metric}), the vector
\[
\nabla f=g^\sharp(\dd f)=\partial_rf\;\partial_r+\frac{1}{r^2}\partial_\varphi f\;\partial_\varphi .
\]
The familiar $\frac1r\partial_\varphi f\,\hat e_\varphi$ is the same thing written in the \emph{unit} basis $\hat e_\varphi=\frac1r\partial_\varphi$.
\end{example}

% REWRITTEN: sphere azimuth phi -> varphi, chart written \phi (docs/NOTATION.md)
\begin{example}[{induced metric on the sphere (from the lecture notes I4)}]
The map $\phi^{-1}(\theta,\varphi)=(r\sin\theta\cos\varphi,\,r\sin\theta\sin\varphi,\,r\cos\theta)$ parametrizes $S^2_r$. Its partial derivatives, \ie the pushforwards of $\partial_\theta,\partial_\varphi$, are
\[
v_\theta=(r\cos\theta\cos\varphi,\ r\cos\theta\sin\varphi,\ -r\sin\theta),\qquad
v_\varphi=(-r\sin\theta\sin\varphi,\ r\sin\theta\cos\varphi,\ 0).
\]
Using the Euclidean dot product of $\R^3$: $|v_\theta|^2=r^2$, $|v_\varphi|^2=r^2\sin^2\theta$, $v_\theta\cdot v_\varphi=0$. Hence the induced metric is
\[
g_{S^2_r}=r^2\big(\dd\theta\otimes\dd\theta+\sin^2\theta\,\dd\varphi\otimes\dd\varphi\big).
\]
This is \emph{the} metric of the sphere, used below to compute Christoffel symbols, curvature and geodesics. For the unit sphere ($r=1$): $g=\mathrm{diag}(1,\sin^2\theta)$, $g^{-1}=\mathrm{diag}(1,1/\sin^2\theta)$.
\end{example}


\begin{synthesis}[pull-back appears three times]
Pull-back $\phi\pb$ is always ``precompose with the map''. (i) For a curve $\gamma$, $\gamma\pb\omega=\omega(\dot\gamma)\,\dd t$ is how a $1$-form is integrated (Section~\ref{sec:int1}); (ii) for a
chart or an embedding it produces the components of a tensor or the induced metric, as in the example above where the Euclidean metric pulled back to the sphere gives $r^2(\dd\theta^2+\sin^2\theta\,\dd\varphi^2)$; (iii) for a parametrised surface
it produces the integrand of a $k$-form (Section~\ref{sec:intk}). It commutes with $\wedge$ and with $\dd$, which is why forms, and only forms, can be integrated without choosing a chart.
\end{synthesis}
